%%%%%%%%%%%%%%%%%%%%%%%%%%%%%%%%%%%%%%%%%
%%
%%          Section: Abstract
%%
%%%%%%%%%%%%%%%%%%%%%%%%%%%%%%%%%%%%%%%%%
% IEEE conference abstracts are a single unstructured paragraph. The ESEM
% structured version (Background / Aims / Method / Results / Conclusions) is
% preserved in the comment block at the bottom for reference.
\begin{abstract}
Classifying issue reports as bugs, feature requests, or questions helps developers triage them. Fine-tuned large language models (LLMs) give the best reported results, but fine-tuning needs a training run, reaches its best scores only with labeled issues pooled across projects, and must be repeated for new labels or a newer model. Retrieval-augmented few-shot prompting (RAG) instead shows the LLM the most similar labeled issues as examples and needs no training. To our knowledge, it has not been systematically evaluated for this task or compared with Low-Rank Adaptation (LoRA) fine-tuning. We conduct an empirical study with four Qwen2.5-Instruct models and 6{,}600 issue reports from eleven projects. \rag, like prior methods, most often labels questions as bugs, so a variant (\frag) removes bug-labeled examples when the retrieved issues suggest a question. Using only the target project's 300 labeled issues, \frag\ comes within 1.0 point of fine-tuning on all 3{,}300 labeled issues of the eleven projects in macro~$F_1$, a difference that is not statistically significant, and needs no training and 33\% less peak GPU memory on average. Without the filter, \rag\ trails fine-tuning by 2.8 points. Given the same labeled issues, both RAG variants outperform fine-tuning at every model size. RAG's inference is slower, and the two approaches make different errors: fine-tuning finds more bugs but over-predicts the bug label. With this filter, retrieval-augmented few-shot prompting preserves the macro~$F_1$ of LoRA fine-tuning with $11\times$ less labeled data per project, no training, and a third less GPU memory.
\end{abstract}

\begin{IEEEkeywords}
Issue classification, Retrieval-augmented generation, Fine-tuning, Large language models, Mining software repositories
\end{IEEEkeywords}

\begin{comment}
%% ESEM 2026 structured abstract (LIPIcs submission), kept for reference.
\myparagraph{Background.} Issue report classification (IRC) accelerates modern software maintenance, yet it remains a challenging task. State-of-the-art automated IRC methods fine-tune Large Language Models (LLMs) with Low-Rank Adaptation (LoRA), which is computationally expensive and must be repeated to incorporate newly labeled issue reports or whenever the model is swapped.

\myparagraph{Aims.} We empirically investigate the hypothesis that low-cost, retrieval-augmented generation (RAG)-based approaches can achieve competitive IRC performance with the LoRA fine-tuning baseline. If validated, this would introduce the possibility of cost-efficient IRC solutions; An avenue that, to our knowledge, remains unexplored in the literature.

\myparagraph{Method.} To test this hypothesis, we conduct a case study on the Qwen2.5-Instruct family, a set of open-source models with strong public-benchmark performance, evaluated at four parameter sizes to analyze the effect of model scale. We propose three retrieval-based methods, \knn, \rag, and \frag, each of which works by retrieving the top-$k$ most similar issues from an index of labeled issues and reasoning on them. Specifically, \knn\ assigns the label via a similarity-weighted vote over the retrieved neighbors; \rag\ presents the neighbors as classified examples in a few-shot LLM prompt; and \frag\ extends \rag\ with a more balanced retrieval scheme that better covers the label space. We evaluate all three methods against a LoRA fine-tuning baseline on an 11-project benchmark, under both \emph{project-specific} and \emph{project-agnostic} data scope settings.

\myparagraph{Results.} With project-specific data, \rag\ outperforms fine-tuning by 0.021--0.040 macro $F_1$ across all model sizes. Fine-tuning on cross-project data beats project-specific \rag\ by only 0.029 aggregate macro $F_1$ despite $11\times$ more labeled data. The balanced retrieval approach, \frag, closes this performance gap, achieving statistical equivalence with fine-tuning when using a retrieval fallback for invalid LLM outputs. Averaged across model sizes, retrieval-augmented few-shot prompting requires 33\% less peak GPU memory than fine-tuning.

\myparagraph{Conclusions.} Under our experimental setup and case study, RAG-based approaches demonstrate competitive performance against LoRA fine-tuning, suggesting they represent a viable and cost-efficient alternative for IRC, one that the literature has yet to fully explore.
\end{comment}
